\chapter{Lezione 10}
\label{chap:lezione_10} 

\begin{flushright}
\textit{Data: 30/10/2025}
\end{flushright}

\section{Modelli di Crescita delle Reti: Il Modello Barabási-Albert}

Si analizzano i modelli dinamici per la generazione di reti, con particolare riferimento al modello Barabási-Albert (BA), introdotto nel 1999 per spiegare l'emergere delle proprietà \textit{Scale-Free} nelle reti reali.

\subsection{Meccanismo di Attaccamento Preferenziale}

Il modello si fonda sul meccanismo dell'\textit{Attaccamento Preferenziale} (Preferential Attachment), riassumibile nel principio "the rich get richer". 
La dinamica di evoluzione della rete è definita come segue:
\begin{enumerate}
    \item \textbf{Condizione Iniziale:} Si parte da un nucleo di $n_0$ nodi esistenti.
    \item \textbf{Crescita:} Ad ogni passo temporale $t$, viene aggiunto un nuovo nodo alla rete.
    \item \textbf{Connessione:} Il nuovo nodo crea $m$ collegamenti ($m \le n_0$) verso nodi già presenti nel sistema.
    \item \textbf{Regola di Attaccamento:} La probabilità $\mathbb{P}(k_i)$ che un nuovo collegamento si attacchi al nodo $i$-esimo dipende linearmente dal grado $k_i$ di tale nodo.
\end{enumerate}

Formalmente, la probabilità di connessione è data da:
\begin{equation}
    \mathbb{P}(k_i) = \frac{k_i}{\sum_j k_j}
\end{equation}
Poiché ogni nuovo nodo porta $m$ nuovi archi, il denominatore rappresenta la somma totale dei gradi nella rete al tempo $t$. 
Notiamo che:
\begin{itemize}
    \item Il numero di nodi cresce linearmente: $n(t) = n_0 + t$.
    \item Il numero totale di archi al tempo $t$ è $m_{tot} = m \cdot t$ (assumendo $t$ grande e trascurando i collegamenti iniziali).
    \item La somma dei gradi è il doppio del numero di archi: $\sum_j k_j = 2 m t$.
\end{itemize}
Per definizione, questo modello genera una rete con una Componente Gigante (GC) unitaria.

\begin{figure}[h]
    \centering
    \begin{tikzpicture}
        % Nodo centrale Hub
        \node[circle, draw=black, fill=colorC!40, minimum size=1.2cm] (hub) at (0,0) {\textbf{Hub}};
        
        % Nodi periferici esistenti
        \foreach \angle in {0, 60, 120, 180, 240, 300} {
            \node[circle, draw=black, fill=colorA!40, minimum size=0.5cm] (n\angle) at (\angle:2cm) {};
            \draw (hub) -- (n\angle);
        }
        
        % Nuovo nodo
        \node[circle, draw=black, fill=colorE!40, minimum size=0.8cm, label=above:New Node] (new) at (4,1) {};
        
        % Frecce di probabilità
        \draw[->, dashed, thick, colorE] (new) -- (hub) node[midway, above, sloped] {High $P$};
        \draw[->, dashed, thin, gray] (new) -- (n0) node[midway, below, sloped] {Low $P$};
    \end{tikzpicture}
    \caption{Rappresentazione schematica dell'attaccamento preferenziale. Il nuovo nodo ha una probabilità maggiore di connettersi all'Hub (grado elevato) rispetto ai nodi periferici.}
\end{figure}

\subsection{Approssimazione nel Continuo}

Per derivare la distribuzione dei gradi, trattiamo $k_i$ come una variabile continua nel tempo. L'evoluzione temporale del grado del nodo $i$ è descritta dall'equazione differenziale:
\begin{equation}
    \frac{\partial k_i}{\partial t} = m \cdot \mathbb{P}(k_i) = m \frac{k_i}{\sum_j k_j}
\end{equation}
Sostituendo la somma dei gradi $\sum_j k_j = 2mt$:
\begin{equation}
    \frac{\partial k_i}{\partial t} = m \frac{k_i}{2mt} = \frac{k_i}{2t}
\end{equation}
Separando le variabili e integrando tra il tempo di nascita del nodo $t_i$ e il tempo corrente $t$:
\begin{equation}
    \int_{k_i(t_i)}^{k_i(t)} \frac{dk_i}{k_i} = \int_{t_i}^{t} \frac{dt'}{2t'}
\end{equation}
La condizione iniziale è che al momento della nascita $t_i$, il nodo ha grado $m$, ovvero $k_i(t_i) = m$.
\begin{equation}
    \ln\left(\frac{k_i(t)}{m}\right) = \frac{1}{2} \ln\left(\frac{t}{t_i}\right)
\end{equation}
Esponenziando entrambi i membri otteniamo la legge di potenza temporale:
\begin{equation}
    k_i(t) = m \left( \frac{t}{t_i} \right)^{\beta} \quad \text{con } \beta = \frac{1}{2}
\end{equation}

\subsection{Derivazione della Distribuzione del Grado $P(k)$}

La probabilità che un nodo abbia un grado $k_i(t)$ inferiore a un valore soglia $k$ è equivalente alla probabilità che il suo tempo di nascita $t_i$ sia successivo a un certo istante (poiché i nodi più vecchi hanno grado maggiore).
Dalla relazione precedente ricaviamo $t_i$:
\begin{equation}
    k = m \left( \frac{t}{t_i} \right)^{1/2} \implies t_i = t \left( \frac{m}{k} \right)^2
\end{equation}
Assumiamo che i tempi di inserimento $t_i$ siano uniformemente distribuiti nell'intervallo $[0, t]$. La funzione di distribuzione cumulativa per $t_i$ è $F(t_i) = t_i / t$.
La probabilità cumulativa per il grado è:
\begin{equation}
    P(k_i(t) < k) = P\left( t_i > \frac{m^2 t}{k^2} \right) = 1 - P\left( t_i \le \frac{m^2 t}{k^2} \right)
\end{equation}
Sostituendo la distribuzione uniforme (assumendo $n_0 \ll t$):
\begin{equation}
    P(k_i(t) < k) = 1 - \frac{1}{t} \left( \frac{m^2 t}{k^2} \right) = 1 - \frac{m^2}{k^2}
\end{equation}
La densità di probabilità $P(k)$ si ottiene derivando la cumulativa rispetto a $k$:
\begin{equation}
    P(k) = \frac{\partial}{\partial k} \left( 1 - \frac{m^2}{k^2} \right) = -(-2) m^2 k^{-3} = 2m^2 k^{-3}
\end{equation}
Si ottiene quindi il celebre risultato che la distribuzione dei gradi decade come una legge di potenza con esponente $\gamma = 3$:
\begin{tcolorbox}[colback=colorC!15, colframe=colorC!70!colorB,  title=\textbf{Decadimento della Distribuzione dei Gradi a Legge di Potenza}]
\begin{equation}
    P(k) \sim k^{-3}
    \label{eq: 10.1}
\end{equation}
\end{tcolorbox}

\subsection{In-Degree Distribution}

Per analizzare più in dettaglio la struttura della rete, è utile distinguere tra i collegamenti uscenti ed entranti. Nel modello Barabási-Albert, ogni nuovo nodo introdotto nella rete genera $m$ collegamenti verso nodi preesistenti. Pertanto, il grado uscente (out-degree) è fisso per ogni nodo e pari a $m$.
La variabilità del grado totale $k$ è determinata interamente dal grado entrante (in-degree), che denotiamo con $q$.

La relazione tra il grado totale $k$ e il grado entrante $q$ è data da:
\begin{equation}
    k = q + m
\end{equation}
Dove $m$ rappresenta il numero di link che il nodo stabilisce al momento della sua nascita (out-degree costante).

Consideriamo la probabilità di attaccamento preferenziale. La probabilità $\mathbb{P}(k_i)$ che un nuovo nodo si connetta al nodo $i$ è proporzionale al grado totale di $i$. Possiamo riscrivere questa probabilità in termini di in-degree $q_i$:
\begin{equation}
    \mathbb{P}(q_i) = \frac{k_i}{\sum_j k_j} = \frac{q_i + m}{\sum_j (q_j + m)}
\end{equation}
Al denominatore, la somma è estesa a tutti i nodi presenti al tempo $t$. Se al tempo $t$ ci sono $t$ nodi (assumendo $n_0 \ll t$), la somma dei gradi totali è pari a due volte il numero totale di archi ($m \cdot t$):
\begin{equation}
    \sum_j k_j = 2 m t
\end{equation}
Pertanto, la probabilità che un singolo nuovo link si connetta a un nodo con in-degree $q$ è:
\begin{equation}
    \mathbb{P}(q) = \frac{q + m}{2 m t}
\end{equation}
Poiché ogni nuovo nodo introduce $m$ nuovi collegamenti, la probabilità complessiva che un nodo con in-degree $q$ riceva un nuovo collegamento durante un passo temporale è:
\begin{equation}
    P(\text{link} \to q) = m \cdot \mathbb{P}(q) = m \frac{q + m}{2 m t} = \frac{q + m}{2 t}
\end{equation}

\subsection{Approccio della Master Equation}

L'approccio della Master Equation ci permette di descrivere l'evoluzione temporale esatta della distribuzione dei gradi.
Definiamo $N(q, t)$ come il numero medio di nodi che possiedono in-degree pari a $q$ al tempo $t$.
La probabilità che un nodo scelto a caso al tempo $t$ abbia in-degree $q$ è data da:
\begin{equation}
    P(q, t) = \frac{N(q, t)}{t}
\end{equation}

L'evoluzione del numero di nodi $N(q, t)$ dal tempo $t$ al tempo $t+1$ è governata dai flussi di probabilità tra gli stati (i diversi valori di $q$).
Dobbiamo considerare tre contributi:
\begin{enumerate}
    \item Guadagno da $q-1$: Un nodo che ha in-degree $q-1$ riceve un nuovo collegamento e il suo grado diventa $q$.
    \item Perdita da $q$: Un nodo che ha già in-degree $q$ riceve un nuovo collegamento e il suo grado diventa $q+1$.
    \item Nuovo nodo: Ad ogni passo temporale viene introdotto un nuovo nodo. Questo nodo nasce con in-degree $q=0$ (poiché i suoi link sono solo uscenti).
\end{enumerate}

\noindent Scriviamo l'equazione di bilancio per $N(q, t+1)$. Per $q > 0$:
\begin{equation}
    N(q, t+1) = N(q, t) + \underbrace{\frac{m[(q-1) + m]}{2mt} N(q-1, t)}_{\text{Guadagno da } q-1} - \underbrace{\frac{m[q + m]}{2mt} N(q, t)}_{\text{Perdita verso } q+1}
\end{equation}
Semplificando i termini $m$ e $t$:
\begin{equation}
    N(q, t+1) = N(q, t) + \frac{q - 1 + m}{2t} N(q-1, t) - \frac{q + m}{2t} N(q, t)
\end{equation}

\noindent Per $q = 0$: il nuovo nodo entra sempre con in-degree 0. Inoltre, nodi con $q=0$ possono ricevere un link e passare a $q=1$.
\begin{equation}
    N(0, t+1) = N(0, t) + 1 - \frac{m}{2t} N(0, t)
\end{equation}
Il termine $+1$ rappresenta il nuovo nodo aggiunto al sistema al tempo $t+1$.

\subsubsection{Soluzione nello Stato Stazionario}

Cerchiamo una soluzione stazionaria in cui la distribuzione di probabilità $P(q)$ non dipenda dal tempo per $t \to \infty$.
Assumiamo $N(q, t) = t P(q)$ e sostituendo nella Master Equation per $q > 0$:
\begin{equation}
    P_q = \frac{1}{2}(q - 1 + m) P_{q-1} - \frac{1}{2}(q + m) P_q
\end{equation}
Raggruppiamo i termini con $P(q)$ a sinistra e otteniamo la relazione ricorsiva fondamentale:
\begin{equation}
    P_q^\infty  = P_{q-1}^\infty  \frac{q + m - 1}{q + m + 2}
\end{equation}

\noindent Analizziamo ora l'equazione per $q=0$ nel limite stazionario. Si trova la seguente relazione:
\begin{equation}
 P_0^\infty = \frac{2}{m + 2}
\end{equation}

\subsection{Soluzione Esatta e Comportamento Asintotico}

Utilizzando la relazione ricorsiva, possiamo scrivere $P(q)$ come prodotto:
\begin{equation}
    P_q^\infty  = P_0^\infty  \prod_{j=1}^{q} \frac{j + m - 1}{j + m + 2}
\end{equation}
Questa espressione può essere riscritta utilizzando la funzione Gamma di Eulero, $\Gamma(x)$, sfruttando la proprietà $\Gamma(x+1) = x \Gamma(x)$.
La relazione ricorsiva $P(q) / P(q-1)$ suggerisce la forma:
\begin{equation}
    P(q) = \frac{2}{m+2} \frac{\Gamma(q+m) \Gamma(m+3)}{\Gamma(m) \Gamma(q+m+3)}
\end{equation}
Verifichiamo i termini costanti. Sapendo che $P(0) = \frac{2}{m+2}$, l'espressione generale diventa:
\begin{equation}
    P(q) = \frac{\Gamma(m+3)}{\Gamma(m)} \frac{\Gamma(q+m)}{\Gamma(q+m+3)}
\end{equation}
Il pre-fattore $\frac{\Gamma(m+3)}{\Gamma(m)} = (m+2)(m+1)m! / (m-1)! \dots$ è una costante di normalizzazione.

\noindent  Per studiare il comportamento asintotico per grandi $q$ ($q \gg 1$), utilizziamo l'approssimazione di Stirling o la proprietà asintotica del rapporto di funzioni Gamma:
\begin{equation}
    \frac{\Gamma(z+a)}{\Gamma(z+b)} \approx z^{a-b} \quad \text{per } z \to \infty
\end{equation}
Nel nostro caso, $z = q$, $a = m$, $b = m+3$.
\begin{equation}
    \frac{\Gamma(q+m)}{\Gamma(q+m+3)} \approx q^{m - (m+3)} = q^{-3}
\end{equation}
Quindi, la distribuzione del grado entrante segue una legge di potenza:
\begin{equation}
    P(q) \sim q^{-3}
\end{equation}

Poiché per grandi $q$ si ha $k \approx q$ (dato che $k = q+m$ e $m$ è costante), la distribuzione del grado totale $P(k)$ avrà lo stesso esponente asintotico:
\begin{equation}
    P(k) \sim k^{-3}
\end{equation}
Questo conferma rigorosamente il risultato riportato nell'Equazione \ref{eq: 10.1}, dimostrando che l'esponente $\gamma = 3$ è una proprietà robusta del modello Barabási-Albert.


\newpage
\section{Dinamiche di Diffusione Epidemica}

La modellizzazione matematica in epidemiologia vanta una storia lunga e ricca, che risale agli anni '20 con la teoria di Kermack-McKendrick. L'idea fondamentale, seppur semplice, si rivela estremamente potente: è possibile suddividere la popolazione in diversi \textit{compartimenti} che rappresentano gli stadi della malattia e utilizzare le dimensioni relative di tali compartimenti per modellare l'evoluzione temporale del contagio.

Sebbene nasca in ambito biologico, questa teoria ha trovato una forte sinergia con la \textit{Network Science}. Oggi, i modelli epidemici sono utilizzati per comprendere non solo la diffusione di patogeni biologici, ma anche la propagazione di idee, virus informatici o "meme" sui social network.
Le domande centrali a cui questi modelli cercano di rispondere sono:
\begin{itemize}
    \item Come si diffonde un'infezione all'interno di una popolazione?
    \item La malattia evolverà in una pandemia globale o si estinguerà localmente?
    \item Qual è la strategia di vaccinazione ottimale?
\end{itemize}

\subsubsection{Dinamica Co-evolutiva}
È fondamentale considerare che esistono due popolazioni interagenti:
\begin{enumerate}
    \item \textbf{Ospiti (Hosts):} Gli individui che contraggono il virus.
    \item \textbf{Patogeni (Virus):} Gli agenti infettivi che evolvono per eludere il sistema immunitario dell'ospite.
\end{enumerate}
Questo sistema è intrinsecamente co-evolutivo. Ad esempio, nel caso dell'influenza stagionale, i vaccini devono essere aggiornati regolarmente per adattarsi ai ceppi dominanti che emergono dalla pressione selettiva.

\hfill

\noindent Iniziamo analizzando i modelli in approssimazione di \textbf{well-mixed population} (popolazione ben mescolata), dove si assume che ogni individuo possa interagire con chiunque altro con la stessa probabilità.

\subsection{Il Modello SI (Susceptible-Infected)}
Il modello più semplice divide la popolazione $N$ in due classi:
\begin{itemize}
    \item \textbf{S (Susceptible):} Individui sani ma suscettibili all'infezione.
    \item \textbf{I (Infected):} Individui infetti in grado di trasmettere la malattia.
\end{itemize}
Sia $\beta$ il tasso di infezione, che dipende dalla contagiosità del patogeno e dalla frequenza dei contatti sociali. La dinamica segue la \textbf{Legge di Azione di Massa}: la variazione degli infetti è proporzionale al prodotto tra il numero di suscettibili e la frazione di infetti presenti.
Le equazioni differenziali che governano il sistema sono:
\begin{align}
    \frac{dI(t)}{dt} &= \beta S(t) \frac{I(t)}{N} \\
    \frac{dS(t)}{dt} &= -\beta S(t) \frac{I(t)}{N}
\end{align}
Con la condizione di normalizzazione $S(t) + I(t) = N$.
Passando alle variabili normalizzate $s(t) = S(t)/N$ e $i(t) = I(t)/N$, otteniamo:
\begin{equation}
    \frac{di}{dt} = \beta s i = \beta (1-i) i
\end{equation}
La soluzione di questa equazione differenziale è la \textbf{funzione logistica}:
\begin{equation}
    i(t) = \frac{i(0) e^{\beta t}}{1 - i(0) + i(0)e^{\beta t}}
\end{equation}
Questo modello prevede che, per $t \to \infty$, l'intera popolazione venga infettata ($i \to 1$). È un modello adatto per malattie incurabili senza recupero, ma limitato per la maggior parte dei casi reali.

\begin{figure}[h]
\centering
\begin{tikzpicture}
    \begin{axis}[
        width=8cm, height=6cm,
        xlabel={$t$}, 
        ymin=0, ymax=1.1, xmin=0, xmax=10,
        axis lines=left,
        grid=major,
        legend pos=outer north east
       ]
    % Curva degli Infetti i(t) - Crescita Logistica
    \addplot[domain=0:10, samples=100, thick, colorC] {1/(1 + 99*exp(-1*x))};
    \addlegendentry{$i(t)$}
    
    % Curva dei Suscettibili s(t) - Decadimento (s = 1 - i)
    \addplot[domain=0:10, samples=100, thick, colorE, dashed] {1 - (1/(1 + 99*exp(-1*x)))};
    \addlegendentry{$s(t)$}
    \end{axis}
\end{tikzpicture}
\caption{Andamento complementare di suscettibili $s(t)$ e infetti $i(t)$ nel modello SI.}
\end{figure}

\subsection{Il Modello SIR (Susceptible-Infected-Recovered)}
Per descrivere malattie da cui si può guarire (o morire), acquisendo immunità, si introduce il compartimento \textbf{R (Recovered)}.
\begin{itemize}
    \item $S \to I$: transizione con tasso $\beta$ (infezione).
    \item $I \to R$: transizione con tasso $\mu$ (guarigione/rimozione).
\end{itemize}
Il sistema di equazioni differenziali è:
\begin{align}
    \frac{ds}{dt} &= -\beta s i \\
    \frac{di}{dt} &= \beta s i - \mu i \\
    \frac{dr}{dt} &= \mu i
\end{align}
Dove $\mu$ è il tasso di recupero, inverso del tempo medio di infezione $\tau \sim 1/\mu$. La probabilità di rimanere infetti per un tempo $\tau$ decade esponenzialmente come $P(\tau) \sim e^{-\mu \tau}$.

\subsubsection{Derivazione della Dimensione Finale dell'Epidemia}
Dividendo la prima equazione per la terza, eliminiamo la dipendenza temporale:
\begin{equation}
    \frac{ds}{dr} = \frac{-\beta s i}{\mu i} = -\frac{\beta}{\mu} s
\end{equation}
Definiamo il \textbf{Basic Reproduction Number} come $R_0 = \frac{\beta}{\mu}$. Integrando l'equazione differenziale tra lo stato iniziale (tempo 0) e un tempo generico $t$:
\begin{equation}
    \int_{s(0)}^{s(t)} \frac{ds'}{s'} = -R_0 \int_{r(0)}^{r(t)} dr' \implies \ln\left(\frac{s(t)}{s(0)}\right) = -R_0 (r(t) - r(0))
\end{equation}
Assumendo che all'inizio nessuno sia guarito ($r(0)=0$) e quasi tutti siano suscettibili ($s(0) \approx 1$), otteniamo:
\begin{equation}
    s(t) = e^{-R_0 r(t)}
\end{equation}
Nel limite asintotico $t \to \infty$, il numero di infetti si annulla ($i_\infty = 0$), quindi vale la conservazione $s_\infty + r_\infty = 1$. Sostituendo $s_\infty = 1 - r_\infty$:
\begin{equation}
    1 - r_\infty = e^{-R_0 r_\infty}
\end{equation}
Otteniamo così un'equazione trascendente per la dimensione totale dell'epidemia $r_\infty$:
\begin{equation}
    r_\infty = 1 - e^{-R_0 r_\infty}
\end{equation}

\begin{figure}[h]
\centering
\begin{tikzpicture}
    \begin{axis}[
        width=8cm, height=6cm,
        xlabel={$r_\infty$},
        xmin=0, xmax=1, ymin=0, ymax=1,
        legend pos=south east,
        axis lines=middle
    ]
    \addplot[domain=0:1, dashed, black] {x};
    \addplot[domain=0:1, thick, colorE] {1 - exp(-2.5*x)}; % R0 > 1
    \addplot[domain=0:1, thick, colorB] {1 - exp(-0.8*x)}; % R0 < 1
    
    \node[colorE] at (axis cs: 0.7, 0.9) {$R_0 > 1$};
    \node[colorB] at (axis cs: 0.3, 0.1) {$R_0 < 1$};
    \end{axis}
\end{tikzpicture}
\caption{Soluzione grafica per $r_\infty$. Per $R_0 > 1$  esiste una soluzione non banale $r_\infty \neq 0$. Per $R_0 < 1$ l'unica soluzione è l'assenza di epidemia ($r_\infty = 0$).}
\end{figure}

Il parametro $R_0$ determina la soglia critica:
\begin{itemize}
    \item Se $R_0 < 1$: L'epidemia non decolla (soluzione $r_\infty = 0$).
    \item Se $R_0 > 1$: Esiste una soluzione positiva $r_\infty > 0$, si verifica un'epidemia macroscopica.
\end{itemize}

\newpage
\subsection{Epidemie su Reti e Percolazione}

Le assunzioni di "well-mixed population" sono spesso eccessivamente semplicistiche. Nella realtà, la struttura dei contatti (la rete sociale, i trasporti aerei, ecc.) gioca un ruolo cruciale.
Passiamo quindi all'analisi delle epidemie su \textbf{Reti Complesse}.

\subsection{Probabilità di Trasmissione e Percolazione di Legame}
Consideriamo un modello in cui ogni individuo infetto rimane tale per un tempo fisso $\tau$. La probabilità di trasmettere l'infezione a un vicino connesso è data da $\beta$ per unità di tempo.
La probabilità totale che l'infezione passi attraverso un arco (link) durante il periodo di infezione $\tau$ è definita come \textbf{trasmissibilità} $\phi$ (o $T$):
\begin{equation}
    \phi = 1 - e^{-\beta \tau}
\end{equation}
Questo ci permette di mappare il problema dinamico dell'epidemia su un problema statico di \textbf{Percolazione di Legame} (\textit{Bond Percolation}).
Ogni arco della rete viene dichiarato "occupato" (ovvero capace di trasmettere l'infezione) con probabilità $\phi$.
\begin{itemize}
    \item Se esiste una \textbf{Componente Gigante} (Giant Component) formata dagli archi occupati, l'epidemia può diffondersi a una frazione macroscopica della popolazione.
    \item La soglia epidemica corrisponde esattamente alla soglia di percolazione $\phi_c$.
\end{itemize}


\subsection{Derivazione della Soglia Critica $\phi_c$}
Vogliamo calcolare la probabilità che un nodo \textbf{non} appartenga alla Componente Gigante (GC) dell'infezione.
Definiamo $u$ come la probabilità che seguendo un arco scelto a caso in una direzione, \textbf{non}  si raggiunga la Componente Gigante. 

Consideriamo un nodo $i$ raggiunto tramite un arco. Affinché questo percorso non conduca alla GC, deve verificarsi una delle seguenti condizioni esclusive:
\begin{enumerate}
    \item L'arco che stiamo percorrendo non è attivo (l'infezione non passa). Questo avviene con probabilità $1 - \phi$.
    \item L'arco è attivo (probabilità $\phi$), ma il nodo $i$ a cui arriviamo non è connesso alla GC tramite nessuno dei suoi \textit{altri} archi uscenti.
\end{enumerate}

\noindent Se il nodo $i$ ha grado totale $k+1$, significa che ha altri $k$ archi oltre a quello di arrivo. La probabilità che nessuno di questi $k$ rami porti alla GC è $u^k$ (assumendo indipendenza locale, valida per reti senza troppi cicli corti).

Dobbiamo mediare questa probabilità sulla distribuzione dei gradi dei nodi che incontriamo navigando la rete. Questa non è la distribuzione di grado standard $P(k)$, ma la \textbf{Excess Degree Distribution} (distribuzione del grado in eccesso), denotata con $q_k$.
La probabilità di arrivare a un nodo con grado $k+1$ è proporzionale a $(k+1)P(k+1)$. Normalizzando:
\begin{equation}
    q_k = \frac{(k+1)P(k+1)}{\langle k \rangle}
\end{equation}
Dove $q_k$ rappresenta la probabilità che un nodo raggiunto abbia $k$ collegamenti ulteriori.

Possiamo ora scrivere l'equazione ricorsiva per $u$:
\begin{equation}
    u = \underbrace{(1 - \phi)}_{\text{Arco inattivo}} + \underbrace{\phi \sum_{k=0}^{\infty} q_k u^k}_{\text{Arco attivo } \cdot \text{ rami morti}}
\end{equation}
Introduciamo la \textbf{Funzione Generatrice} $g_1(u)$ per la distribuzione di eccesso:
\begin{equation}
    g_1(u) = \sum_{k=0}^{\infty} q_k u^k
\end{equation}
L'equazione di auto-consistenza diventa quindi:
\begin{equation} \label{eq:self_consistency}
    u = 1 - \phi + \phi g_1(u)
\end{equation}

\subsubsection{Soluzione Grafica e Transizione di Fase}

Cerchiamo le soluzioni dell'equazione (\ref{eq:self_consistency}) nell'intervallo $u \in [0, 1]$.




\begin{itemize}
    \item Una soluzione banale è sempre $u=1$. Infatti, poiché $\sum q_k = 1$, si ha $g_1(1)~=~1$, quindi $1 = 1 - \phi + \phi(1)$ è verificata. Questo corrisponde all'assenza di una Componente Gigante (l'epidemia si estingue).
    \item Una soluzione fisica $u < 1$ (che implica l'esistenza della GC) emerge solo se la curva interseca la bisettrice $y=u$ con una pendenza sufficiente.
\end{itemize}

La transizione di fase (il punto critico) avviene quando la derivata del membro di destra rispetto a $u$, calcolata in $u=1$, è esattamente uguale a 1 (tangenza):
\begin{equation}
    \frac{d}{du} \Big( 1 - \phi + \phi g_1(u) \Big) \Big|_{u=1} = 1
\end{equation}
Calcolando la derivata:
\begin{equation}
    0 - 0 + \phi \cdot g_1'(u) \Big|_{u=1} = 1 \implies \phi_c g_1'(1) = 1
\end{equation}
La soglia critica è dunque:
\begin{equation}
    \phi_c = \frac{1}{g_1'(1)}
\end{equation}

\subsubsection{Calcolo Esplicito dei Momenti}

Dobbiamo calcolare il valore di $g_1'(1)$. Dalla definizione di funzione generatrice:
\begin{equation}
    g_1'(u) = \frac{d}{du} \sum_{k=0}^{\infty} q_k u^k = \sum_{k=1}^{\infty} k q_k u^{k-1}
\end{equation}
Valutando in $u=1$:
\begin{equation}
    g_1'(1) = \sum_{k=0}^{\infty} k q_k
\end{equation}
Questo rappresenta il valor medio della distribuzione di eccesso. Sostituiamo ora la definizione di $q_k = \frac{(k+1)P(k+1)}{\langle k \rangle}$:
\begin{equation}
    g_1'(1) = \sum_{k=0}^{\infty} k \frac{(k+1)P(k+1)}{\langle k \rangle}
\end{equation}
Per risolvere questa sommatoria, effettuiamo un cambio di variabile. Poniamo $j = k+1$, da cui $k = j-1$. Quando $k=0$, $j=1$.
\begin{equation}
    g_1'(1) = \frac{1}{\langle k \rangle} \sum_{j=1}^{\infty} (j-1) j P(j)
\end{equation}
Espandiamo il prodotto $(j-1)j = j^2 - j$:
\begin{equation}
    g_1'(1) = \frac{1}{\langle k \rangle} \left( \sum_{j} j^2 P(j) - \sum_{j} j P(j) \right)
\end{equation}
Riconosciamo i momenti della distribuzione originale $P(k)$: $\sum j^2 P(j) = \langle k^2 \rangle$ e $\sum j P(j) = \langle k \rangle$.
Quindi:
\begin{equation}
    g_1'(1) = \frac{\langle k^2 \rangle - \langle k \rangle}{\langle k \rangle}
\end{equation}

\subsubsection{Risultato Finale per la Soglia Epidemica}

Sostituendo l'espressione di $g_1'(1)$ nella formula della soglia critica $\phi_c = 1/g_1'(1)$, otteniamo il risultato fondamentale:

\begin{tcolorbox}[colback=colorC!15, colframe=colorC!70!colorB,  title=\textbf{Soglia Epidemica}]
\begin{equation} \label{eq:threshold}
    \phi_c = \frac{\langle k \rangle}{\langle k^2 \rangle - \langle k \rangle}
\end{equation}
\end{tcolorbox}


Questa equazione lega la dinamica epidemica (rappresentata da $\phi_c$) esclusivamente alle proprietà topologiche della rete (i momenti $\langle k \rangle$ e $\langle k^2 \rangle$).

\subsection{Implicazioni per la Robustezza delle Reti}

Analizziamo il comportamento della soglia (\ref{eq:threshold}) per diverse topologie di rete.

\subsubsection{Caso 1: Reti Omogenee (Erdős-Rényi)}
In una rete random con distribuzione di Poisson, la varianza $\sigma^2$ è uguale alla media $\langle k \rangle$.
Ricordando che $\sigma^2 = \langle k^2 \rangle - \langle k \rangle^2$, possiamo scrivere:
\begin{equation}
    \langle k^2 \rangle = \langle k \rangle^2 + \sigma^2 = \langle k \rangle^2 + \langle k \rangle
\end{equation}
Sostituendo nel denominatore della soglia critica:
\begin{equation}
    \phi_c = \frac{\langle k \rangle}{(\langle k \rangle^2 + \langle k \rangle) - \langle k \rangle} = \frac{\langle k \rangle}{\langle k \rangle^2} = \frac{1}{\langle k \rangle}
\end{equation}
Se la rete ha grado medio $\langle k \rangle = 4$, allora $\phi_c = 1/4 = 0.25$.
Ciò implica una notevole \textbf{robustezza}: affinché l'epidemia si propaghi (o affinché la rete rimanga connessa sotto attacco casuale), la trasmissibilità deve essere superiore al 25\%. Inversamente, possiamo rimuovere fino al 75\% dei nodi prima che la componente gigante venga distrutta.

\subsubsection{Caso 2: Reti Scale-Free (Barabási-Albert)}
Le reti reali (Internet, WWW, reti sociali) seguono una legge di potenza $P(k) \sim k^{-\gamma}$ con $2 < \gamma \le 3$.
Una proprietà peculiare di queste distribuzioni è che, nel limite di rete infinita ($N \to \infty$), il momento secondo diverge:
\begin{equation}
    \langle k^2 \rangle \to \infty
\end{equation}
Se il denominatore nella (\ref{eq:threshold}) tende a infinito, la soglia critica tende a zero:
\begin{equation}
    \phi_c = \frac{\langle k \rangle}{\infty} \longrightarrow 0
\end{equation}

Questo risultato dimostra l'estrema \textbf{fragilità} delle reti scale-free rispetto alla diffusione virale.
\begin{itemize}
    \item Poiché $\phi_c = 0$, qualsiasi virus anche con trasmissibilità infinitesimale è in grado di diffondersi e diventare pandemico.
    \item Non esiste una soglia di sicurezza naturale.
    \item Un numero infinitesimale di nodi iniziali è sufficiente a sostenere il contagio 
\end{itemize}

